ファイルへの書き込み権限がないため、修正済みLaTeX全文をここに出力します。

\section{違いを検証する - 統計学的検定 -}
  \subsection{【復習】統計量}
    以前「データを要約する」の章で, データの特徴を数値で表すという考え方を学んだ.
    たとえば交通費の記録があっても, 一つひとつを眺めるだけでは全体の傾向はわからないが,
    平均値などで要約すれば, どのくらい使ったかがつかみやすくなる.

    このような, データの特徴を表す数値を\textbf{統計量}と呼ぶのだった.

    統計量はデータを要約したものであるから, データどうしを比較したり差を調べたりするときに役に立つ.
    このあと学ぶ統計的検定では, 統計量を使ってデータに差があるかどうかを調べていく.

  \subsection{統計的に有意とは？}

    たとえば, 新しい栄養補給サプリを試したとき,
    「これは本当に効いているのだろうか」と考えたことはないだろうか.
    あるいは, 「○○県出身の人は△△の傾向が強い」というような主張を耳にしたとき,
    「そんなことはないぞ」と感じた経験があるかもしれない.
    このような問題に対して,
    「本当に効果があるのか（△△の傾向に差があるのか）, 単なる偶然なのか」
    を判断するための考え方のひとつが,
    \textbf{統計的有意性}である.

    統計的有意性の考え方では, 手元にあるデータが「普通の」データに対してどのくらい珍しいか, を考える.
    例えば, 飲めばテストの成績があがる薬があったとしよう.
    この薬に効果があるかどうかを考えるには, 薬を飲んでいなかった場合の成績と薬を飲んだ場合の成績を比べればよさそうである.
    そして, 成績の差は, もし薬に効果がなければ, 薬を飲んでいなかった場合と比較してほとんど差がないはずだ.
    逆に大きく差があれば, （もしかしたら逆効果かもしれないが）薬になんらかの効果がある, といえそうだ.
    このように今あるデータを統計的に意味があるかどうか調べることを\textbf{統計的検定}という.
    統計的検定が言おうとする「有意」とは, 手元のデータと普通のデータの違いが, 「ふつう起こり得るばらつきの幅」をはっきり越えている, という意味である.
    やることは単純だ.
    まず, 手元のデータと普通のデータの比べ方を決めてひとつの統計量として表す.
    つぎに, その統計量が「差がない」場合の統計量の分布においてどのくらい珍しいかを調べる.
    その珍しさが「よくある」ような幅に収まるなら「偶然で説明できる」と判断し, 明らかにはみ出していれば「偶然では説明しにくい」と判断して, \textbf{統計的に有意}と言うのだ.

    \noindent\textbf{用語}
    \begin{itemize}
      \item \textbf{帰無仮説}：差はない／効果はゼロ, という基準となる仮説.
      \item \textbf{対立仮説}：差がある／効果がある, という主張の仮説.
    \end{itemize}

 検定は, 「差がない世界」をいったん仮定し, その世界でも今見た差がどれくらい起こりうるかを調べ, 偶然では説明しにくければ「差がありそうだ」と判断する方法である.

  たとえば, 新しいサプリを試したとき, 「本当に効いているのか」が知りたいとする.
  帰無仮説として\textbf{効果はゼロ}だと\underline{仮に}考える.
  この世界では, サプリを飲んだグループとそうでないグループは\textbf{統計的には区別がつかない}はずだ.
  このとき, サプリの効果を測る点数の「グループごとの平均点」の差は, 中心がほぼ $0$ で, 広がりは\textbf{データのばらつき}と\textbf{人数}で決まるばらつき方をする（中心極限定理の帰結である）.
  この\textbf{差がない世界での分布}を基準として, 今観測した統計量と同じかそれ以上に極端な値が出る割合を求める.
  この割合を\textbf{p値}と呼ぶ.
  p値が十分小さければ, 「偶然だけでは説明しにくい」と判断し, \textbf{統計的に有意}と述べる. なお, p値は帰無仮説が正しい確率ではなく, 帰無仮説のもとでの「珍しさ」の度合いを表す指標である.

    \begin{enumerate}
      \item \textbf{何を比較するかを決める} \\
        調べたい問題に応じて, 比較すべき特徴が決まる.
        たとえば, 二つのグループの平均値を比べるか,
        データのばらつき（分散）を比べるかといった違いである.
        現在よく知られている統計量の中から,
        問題に適したものを選ぶのが一般的である\footnote{
          ここで選んだ統計量は検定に用いられることから, 検定統計量と呼ばれることもある.
          新たな統計量を設計することも理論上は可能であるが,
          妥当性の検討や分布の導出が必要となり, 現実的ではないことが多い.}.

      \item \textbf{帰無仮説を立てる} \\
        「差がない」「効果がない」とする仮説を立てる.
        「差がゼロ」は数値として基準が明確であるため,
        観察されたデータがそこからどれだけ離れているかを確率的に評価しやすい.
        また, 検定統計量が帰無仮説のもとで従う分布はよく知られており,
        確率の計算に使える.
        この点は, 次の章で具体的な検定の計算を通して確認する.

      \item \textbf{データから統計量を計算する} \\
        実際に得られたデータを使って, 選んだ統計量を計算する.
        たとえば, 二つのグループの平均値の差を求めるといったことを行う.

      \item \textbf{統計量の基準分布をもとに確率を調べる} \\
        「差がない」と仮定したとき, 検定統計量が従う\textbf{基準とする分布（帰無仮説のもとでの分布）}を用いて,
        観察された統計量と同じかそれ以上に極端な結果が得られる確率を求める.
        この確率が\textbf{p値}である.
        p値がいくつ以下なら有意とするかの基準（\textbf{有意水準}）は, 検定の前に決めておく.

      \item \textbf{p値をもとに判断する} \\
        有意水準よりもp値が小さければ,
        帰無仮説は成り立ちがたいと判断し, 対立仮説の方がより現実に合っていると考える,
        すなわち「統計的に意味のある差がある」と判断する.
        そうでなければ, 「帰無仮説は棄却できない（帰無仮説を捨てることはできない）」とみなす.

        なお, 検定では「差があるか（両側検定）」なのか,
        「特定の方向に差があるか（片側検定）」なのかを,
        問題の目的に応じて選ぶ必要がある.
    \end{enumerate}


    ただし, p値が小さくなかったからといって, 「差がまったくない」とは言いきれない.
    たまたま今回のデータでは明確な証拠が見つからなかっただけかもしれないからだ.
    「差がないこと」を直接示すには, 「同等性検定」という別の手順が必要になる.
    本書では扱わないが, 必要になったときに調べればよい.


  \subsection{使われる統計量と検定の種類}

    統計的検定では, データから計算された統計量をもとに,
    帰無仮説が成り立つかどうかを調べると述べた.
    どの統計量を使うかは, 問いの内容やデータの種類によって決まり, たくさんの種類がある.

    以下に代表的な統計量と, それに基づく検定の関係を示す：

    \begin{itemize}
      \item \textbf{t統計量}：\quad 二つのグループの「平均値」の差を評価する.
        この統計量を用いた検定が, いわゆる「t検定」である.

      \item \textbf{カイ二乗統計量}：\quad 観測されたカテゴリの「割合」や「傾向」が,
        期待される割合と比べてどの程度偏っているかを評価する.
        この統計量を用いた検定が「カイ二乗検定」である.

      \item \textbf{F統計量}：\quad 二つ以上のグループの「ばらつき（分散）」の差を評価する.
        この統計量を用いた検定が「F検定」や「分散分析（ANOVA）」である.
    \end{itemize}

    これらの統計量はいずれも,
    「差がない」と仮定したときに従うとされる分布（t分布, カイ二乗分布など）が知られており,
    実際に得られた統計量がその分布の中でどれくらい珍しいかを調べることで,
    意味のある差かどうかを判断する.

    統計的検定では,
    問いに応じて適切な統計量を選び,
    それに基づいてデータを評価する.


  \subsection{具体的な検定の例：語学学習アプリの効果}
    ある語学学習アプリが新たに開発され, その効果を検証したいとする.
    開発者は, 「このアプリを使えば英単語の習得が効率よく進み, テストの得点が上がるはずだ」と考えている.

    そこで, 学習者を無作為に二つのグループに分け,
    一方にはアプリを使ってもらい（\textbf{アプリ使用群}）,
    他方には従来通りの方法で学習してもらった（\textbf{通常学習群}）.
    一定期間学習した後で, 同じ英単語テストを受けてもらい, 平均点に差があるかどうかを比較する.

    このように, 別々の人を対象にした二つのグループを比べる方法は,
    \textbf{対応のない二群の比較}と呼ばれる.

    \begin{enumerate}
      \item \textbf{何を比較するかを決める} \\
        今回の目的は, アプリを使ったグループと使っていないグループの\textbf{平均点}を比較することである.
        よって, 検定統計量としては二つの平均値の差に基づく\textbf{t統計量}を用いるのが適切である.

      \item \textbf{帰無仮説を立てる} \\
        「アプリの有無によって英単語テストの平均点に差はない」と仮定する.
        これが\textbf{帰無仮説}である.
        一方で, 「アプリ使用群の方が平均点が高い」という主張は\textbf{対立仮説}となる.

        なお, 今回は「アプリ使用群の方が高いか」を検証しており,
        これは特定の方向への差に注目する\textbf{片側検定}にあたる.
        一方で, 「単に差があるかどうか」を検討したい場合は,
        \textbf{両側検定}を行う.
        検定の方向は, 問題の目的や仮説の立て方によってあらかじめ定めておく必要がある.

      \item \textbf{データから統計量を計算する} \\
        実際に二つのグループの平均値とばらつきを調べ,
        それらを使って\textbf{t統計量}と呼ばれる数値を計算する.
        このt統計量がどのように計算されるかは, 次の章で詳しく見ていく.

      \item \textbf{t分布をもとにp値を求める} \\
        「平均点に差がない」と仮定したときに,
        得られたt統計量と同じかそれ以上に極端な値が出る確率を,
        t分布を使って計算する.
        これが\textbf{p値}である.

      \item \textbf{p値をもとに判断する} \\
        あらかじめ「p値が0.05未満なら有意」と決めていたとしよう.
        このとき, 実際に求めたp値が0.05より小さければ,
        帰無仮説は成り立ちがたいと判断し,
        「アプリ使用群の方が平均点が高い」という主張を支持する.
    \end{enumerate}

    ところで,
    「同じ人の学習前と学習後を比べたほうが確実ではないか」と感じた人もいるかもしれない.

    実際, 一人の人にアプリを使う前と使った後でテストを受けてもらい,
    点数の変化を見るという方法もある.
    このように, 同じ対象から得られた二つの値を比較する方法は,
    \textbf{対応のある検定}と呼ばれる.

    このように検定の方法には複数の選択肢があり,
    どれを選ぶかは, データの集め方や分析の目的によって変わる.
    一つの検定結果だけでなく, 効果の大きさや実験デザインも合わせて判断する必要がある.


  \subsection{具体的な検定の例：血液型と性格に関係はあるか？}

    一昔前に, 血液型によって性格がわかる, という診断が流行ったことがある.
    「A型は几帳面, B型はマイペース」といった話を聞いたことがある人も多いだろう.
    こうした主張には科学的根拠が疑わしいものも多いが,
    「血液型と性格に関係があるのか」という問いを統計的に扱うとどうなるか.

    このように, それぞれ「いくつかのグループに分けられるデータ」どうしの関係を調べるときに使える手法のひとつに
    \textbf{カイ二乗検定}がある.

    たとえば以下のように, 100人に「血液型」と「性格タイプ」についてアンケートをとったとする.
    性格については, 見やすさのために「几帳面」と「大雑把」の二つに絞って分類している：

    \begin{center}
      \begin{tabular}{l|cccc|c}
             & A型 & B型 & O型 & AB型 & 合計 \\
      \hline
      几帳面 & 18  & 6   & 10  & 8    & 42 \\
      大雑把 & 7   & 19  & 15  & 17   & 58 \\
      \hline
      合計   & 25  & 25  & 25  & 25   & 100
      \end{tabular}
    \end{center}

    この表では, 各血液型と性格タイプの組み合わせごとに人数が記録されている.
    たとえば A型の中で「几帳面」と答えた人は18人, 「大雑把」と答えた人は7人である.

    このような表（\textbf{クロス集計表}あるいは\textbf{分割表}と呼ばれる）をもとに,
    「血液型と性格タイプは関係があるのか」を調べるには, 以下の手順をとる.

    \begin{enumerate}
      \item \textbf{帰無仮説を立てる} \\
        「血液型と性格タイプは無関係である」と仮定する.
        つまり, 血液型によって性格タイプの割合が変わることはない, というのが帰無仮説である.

      \item \textbf{期待される人数を計算する} \\
        血液型と性格が本当に無関係なら,
        各セルに入るべき人数は, 行と列の合計から割り算で求められる.
        たとえば A型で几帳面な人は $25 \times \frac{42}{100} = 10.5$ 人程度と期待される.
        このような「期待される人数」を, 表のすべてのセルについて計算する.

      \item \textbf{カイ二乗統計量を計算する} \\
        実際の人数と, 無関係だった場合に期待される人数とのズレを調べる.
        ズレが大きいほど, 「無関係とは言いにくい」と判断される.
        このズレをもとに, カイ二乗検定では\textbf{カイ二乗統計量}と呼ばれる数値を計算する.
        具体的な計算式は, 次の章で紹介する.

      \item \textbf{カイ二乗分布をもとにp値を求める} \\
        帰無仮説が正しいとしたとき, 得られた $\chi^2$ の値がどの程度まれかを,
        カイ二乗分布を使って調べる.
        ここで求めた確率が\textbf{p値}である.

      \item \textbf{p値をもとに判断する} \\
        あらかじめ「p値が0.05未満なら有意」と決めていたとすると,
        実際のp値がこれより小さければ,
        「血液型と性格は無関係とは考えにくい」と判断する.
    \end{enumerate}

    なお, 実際にはこうした調査では,
    「血液型を意識して回答する人がいる」,
    「性格の分類が恣意的である」など,
    さまざまなバイアスや問題点がある.
    検定結果だけで「血液型と性格に関係がある」と結論づけるのは危険であり,
    調査設計やデータの集め方も含めて慎重に判断する必要がある.
    カイ二乗検定は, このようなカテゴリデータどうしの関係を調べるときに使える手法である.


  実際のデータ分析では, 検定の結果をどう解釈するか,
  検定を行う際にどのような注意が必要かという問題も出てくる.

  次の章では,
  検定の数学的な計算方法にもう少し踏み込み,
  p値の意味や, 検定に潜む落とし穴, 効果量やサンプルサイズとの関係についても扱う.
